\chapter{The diff loop in depth}
\label{ch:diffloop}

\begin{aside}
Most TUI libraries spend the bulk of a frame in their diff. Ten
microseconds shaved here is ten microseconds your application can
spend thinking. \maya{}'s diff is two pages of code; this chapter
is twenty.
\end{aside}

Chapter~\ref{ch:simd} introduced the vector instructions. This chapter
zooms in on the loop structure that wraps them: how runs are
detected, how spans become ANSI, what happens at the buffer edges,
and where the loop intentionally trades a little CPU for a lot of
output bandwidth.

\section{The naive baseline}

The simplest correct diff is one cell, one ANSI burst:

\begin{lstlisting}
for (int y = 0; y < h; ++y)
  for (int x = 0; x < w; ++x)
    if (front[y][x] != back[y][x]) {
      emit_cursor_move(x, y);
      emit_sgr(back[y][x].style);
      emit_glyph(back[y][x].cp);
    }
\end{lstlisting}

This is correct but terrible. For an 80x24 grid filled with random
changes you emit roughly 50 bytes of ANSI per changed cell, most of
it cursor moves and SGR sets. \maya{} aims for under 5 bytes per
cell in steady state.

\section{Run detection}

The first optimisation is to detect \emph{runs} of consecutive
changed cells with the same style. Within a run, no cursor move
and no SGR is needed:

\begin{lstlisting}
for (int y = 0; y < h; ++y) {
  int x = 0;
  while (x < w) {
    if (front(x,y) == back(x,y)) { ++x; continue; }
    int start = x;
    Style s = style_of(back(x,y));
    while (x < w
        && front(x,y) != back(x,y)
        && style_of(back(x,y)) == s) ++x;
    emit_run(start, y, s, back, x - start);
  }
}
\end{lstlisting}

Now a 20-cell solid-colour change costs one cursor move, one SGR
set, and 20 UTF-8 glyph encodings: roughly 35 bytes total.

\section{SIMD-accelerated skip}

The \code{while (front(x,y) == back(x,y))} loop is what SIMD
accelerates. \maya{} reads 4 cells (256 bits) at a time on AVX2 or
8 cells on AVX-512:

\begin{lstlisting}
__m256i f = _mm256_loadu_si256(front_row + x);
__m256i b = _mm256_loadu_si256(back_row + x);
__m256i eq = _mm256_cmpeq_epi64(f, b);
uint32_t mask = _mm256_movemask_epi8(eq);
if (mask == 0xFFFFFFFFu) { x += 4; continue; }
// fall back to scalar to find the exact change boundary
\end{lstlisting}

The fast path skips identical 32-byte chunks at memory bandwidth.
The slow path handles boundaries one cell at a time. The
combination is roughly 10x the throughput of a pure scalar loop on
typical TUI screens, which tend to have long unchanged horizontal
spans (margins, padding, backgrounds).

\section{The style cache}

Within a frame, \maya{} caches the last SGR set so it can skip
redundant emissions:

\begin{lstlisting}
Style last = Style::invalid();
for (... each run ...) {
  Style s = run.style;
  if (s != last) { emit_sgr(s); last = s; }
  emit_glyphs(run.cells);
}
\end{lstlisting}

This matters when many small runs share style --- a fullscreen
fill with one accent colour, for example. Without the cache, every
run resets the SGR.

\section{Cursor move optimisation}

Cursor moves come in three flavours:

\begin{itemize}[leftmargin=*]
  \item \textbf{Absolute.} \code{ESC [ y ; x H}, ~8 bytes.
  \item \textbf{Carriage return + down.} \code{\char`\\r\char`\\n},
    2 bytes, only valid for column 0.
  \item \textbf{Implicit.} If the previous emission left the
    cursor exactly where the next run starts, no move is needed.
\end{itemize}

\maya{} tracks the predicted cursor position after each glyph run
and picks the cheapest move. For a vertical column of changes
(typing in a centred input field), the saving over naive absolute
moves is 6 bytes per line.

\begin{tryit}
Set the environment variable \code{MAYA\_TRACE\_ANSI=1} (if your
build was compiled with tracing) and watch a frame. You will see
many runs end with no cursor move at all --- the previous run
left it in the right place.
\end{tryit}

\section{Wide cells and combining marks}

Wide cells (CJK, emoji) take two grid columns but emit as one
glyph. The cell at $x+1$ has width 0 and code point 0; the diff
must treat the pair as one unit, otherwise it would emit half a
glyph and corrupt the screen.

\maya{}'s solution: every cell carries its width in the packed
representation, and the diff loop, on detecting a width-2 cell,
extends the run by one column without emitting anything for the
trailing slot.

\begin{warning}
If you write a custom canvas writer you must zero the trailing
slot of a wide cell. \maya{}'s built-in \code{put} does this.
Forgetting leaves stale glyph data that confuses the diff into
emitting partial sequences.
\end{warning}

\section{Hyperlinks}

OSC-8 hyperlinks are a special case. The link is part of the
style, but the SGR set does not carry it. \maya{} emits a
separate \code{ESC ] 8 ; ; URL \char`\\7} before runs whose link
changes and \code{ESC ] 8 ; ; \char`\\7} to terminate them.

Like SGR, the link is cached: if the next run has the same URL,
no reset.

\section{Bracketed paste and mouse mode}

These are not diffed --- they are emitted once at startup and
once at shutdown, and on any state toggle (entering a modal that
disables mouse, for example). They live outside the buffer.

\section{The emit buffer}

ANSI bytes do not go straight to \code{write(1)}. \maya{}
accumulates them into a \code{std::string} and flushes once at
end of frame:

\begin{lstlisting}
std::string out;
out.reserve(initial_estimate);
diff(front, back, out);
::write(1, out.data(), out.size());
\end{lstlisting}

One write means one syscall per frame. Without batching, every
SGR set and cursor move would be its own syscall (terminals are
line-buffered, but only by convention; raw mode is byte-buffered).

\begin{idea}
\textbf{Syscall amortisation.} Bundling many small writes into
one big write turns N syscalls into one. The kernel cost per
byte goes to zero in the limit.
\end{idea}

\section{When the diff is wrong}

The diff is wrong if the front buffer does not match what is
actually on screen. This happens when:

\begin{itemize}[leftmargin=*]
  \item Another program wrote to the terminal (a stray
    \code{printf}).
  \item The user pressed \code{Ctrl-L} expecting a redraw.
  \item The terminal was resized without \maya{} noticing
    (rare; SIGWINCH catches this).
\end{itemize}

The recovery is a sentinel-flood (see Chapter~\ref{ch:buffers}). The
next frame redraws everything.

\section{Synchronised output}

Some terminals (Kitty, Wezterm, foot) support DEC mode 2026,
which brackets a frame with \code{ESC [ ? 2026 h} and
\code{ESC [ ? 2026 l}. The terminal buffers all the ANSI
between the brackets and applies it atomically. The user never
sees a half-painted frame.

\maya{} probes for this at startup and uses it when available.
Without it, fast-typing users on slow terminals may see
mid-frame flicker.

\begin{funfact}
DEC mode 2026 was specified by the Contour terminal author
Christian Parpart in 2020 and ratified by xterm a year later.
It is the first new private mode to gain broad adoption since
mouse tracking.
\end{funfact}

\section{Recap}

\begin{recap}
\begin{itemize}[leftmargin=*]
  \item Naive diff is correct but bandwidth-expensive.
  \item Run detection cuts SGR and cursor overhead.
  \item SIMD accelerates the skip over unchanged regions.
  \item Style and link caches avoid redundant emissions.
  \item Cursor move optimisation picks the shortest sequence.
  \item Wide cells are one unit; the trailing slot is silent.
  \item DEC 2026 prevents mid-frame flicker on capable
    terminals.
\end{itemize}
\end{recap}

\section*{Workshop}

\begin{enumerate}[leftmargin=*]
  \item Instrument the diff to count: total cells, changed
    cells, runs emitted, bytes emitted. Display these in a
    status line.
  \item Force a sentinel flood on every keystroke. Measure the
    wire bandwidth increase --- it should be 100--1000x.
  \item Detect at startup whether the terminal supports DEC
    2026 by sending a query and parsing the response. Skip the
    bracketing if it does not.
\end{enumerate}
